\section{Scope and Limitations}\label{sec:limits}

This section is treated as a structural component of the report,
not as a concession. It records, in plain language, what the
present version covers and what it does not, so that subsequent
versions can extend the scope without ambiguity.

\paragraph{V1 covers.}
Stage~3 partial closure: items 1, 2, 3, 4a and 4b-dirent are
\textsc{closed} in the source tree and pushed to the public
repository (\cref{sec:implementation}). The on-disk format
v3 is stable. The universal protection scheme
\texttt{INODE\_UNIVERSAL} is the active scheme. Static and
functional verification (checkpatch, four xfstests, dmesg sweep,
module hash uniformity) are reported in \cref{sec:evaluation}.
The I/O baseline of \cref{tab:iobench} is the reference
against which v2 will measure regression after the format~v4
transition.

\paragraph{V1 does not cover.}
\begin{itemize}
\item Item~4 (Shannon entropy field, \texttt{struct ftrfs\_rs\_event}
  growth from 24 to 40 bytes, superblock layout migration to 13~RS
  subblocks). The design is frozen in
  \texttt{Documentation/roadmap.md} but no kernel or user-space
  code has been written. v2 will report it.
\item User-space CVE remediation. The hardened build reports
  239~unpatched CVEs distributed across user-space packages
  (\texttt{openssl}, \texttt{glibc}, \texttt{libxml2},
  \texttt{libpam}, \texttt{python3}, etc.). No backport patches
  or version upgrades have been applied; this is the scope of
  the v0.4 milestone.
\item Stage~4 (universal data-block protection). The three
  reserved scheme values \texttt{UNIVERSAL\_INLINE},
  \texttt{UNIVERSAL\_SHADOW} and \texttt{UNIVERSAL\_EXTENT} are
  defined in the header but no implementation backs any of them.
\item Stage~5 (offensive security review). No syzkaller campaign,
  KASAN-instrumented stress run, or formal adversarial review has
  been conducted.
\item Hardware-on-physical-platform measurement. The bench cluster
  is QEMU/TCG arm64; no measurements on physical embedded hardware
  are reported.
\item Comparative benchmarking against
  \texttt{ext4 + dm-integrity + dm-verity} on equivalent hardware.
\end{itemize}

\paragraph{Methodological limits.}
The author is the sole contributor; no peer review beyond the
public mailing-list interactions on linux-fsdevel has shaped this
text. The bench harness has known limitations: \texttt{dd}-based
timing at millisecond resolution produces quantization artifacts
(\cref{sec:evaluation}); host load isolation is not enforced
during measurement. These limitations are documented but not
remediated in v1.

The acknowledgement of these gaps is itself an application of
axis~3.2 (MIL-STD-882E hazard tracking practice): rather than
present the project as more complete than it is, we record the
state, the gap, and the resolution path.
